% =============================================================
% latex/ean13-tikz.tex
%
% TikZ-Renderer für EAN-13-Strichcodes.
%
%   \eanbarcodebits{<95 Bits, kommasepariert>}
%                  {<Erstziffer>}{<linke 6>}{<rechte 6>}
%
% Modulbreite ist 1mm (über die tikzpicture-x/y-Einheit skalierbar),
% Strichhöhe 30 Module für Daten und 33 für Guards. Klarschrift wird
% darunter mitgesetzt — Erstziffer in der linken Hellzone, die zwei
% Sechsergruppen je unter „ihrer" Hälfte.
%
% Convenience-Wrapper für die im Buch verwendeten EANs sind unten
% definiert (\barcodehonig, später ggf. weitere).
%
% Voraussetzungen: TikZ (über praeambel.tex geladen).
% =============================================================

\newcommand{\eanbarcodebits}[4]{%
  \begin{tikzpicture}[x=1mm, y=1mm, baseline=(current bounding box.south)]
    % Alle Striche enden oben auf y=33. Daten-Striche starten unten bei
    % y=0 (auf der Baseline), Guard-Striche reichen 3 Module weiter
    % nach unten bis y=-3 — auf gleicher Höhe wie die Klarschrift, aber
    % lateral zwischen ihren Sechsergruppen.
    \foreach \bit [count=\j from 0] in {#1}{%
      \ifnum\bit=1
        \def\bo{0}%
        \ifnum\j<3\relax\def\bo{-3}\fi
        \ifnum\j>91\relax\def\bo{-3}\fi
        \ifnum\j>44\relax\ifnum\j<50\relax\def\bo{-3}\fi\fi
        \fill[black] (\j,\bo) rectangle ({\j+1},33);%
      \fi
    }%
    % Klarschrift in Source Code Pro (OCR-B-Tradition). Ziffern stehen
    % auf gleicher Höhe wie die Guard-Strich-Enden, aber lateral
    % versetzt: Erstziffer in der Hellzone, dann zwei Sechsergruppen
    % unter ihrer Hälfte, mit Lücken über den Start-/Mittel-/End-Guards.
    \node[anchor=north, font=\ttfamily\large] at (-5,-0.5) {#2};%
    \node[anchor=north, font=\ttfamily\large] at (24,-0.5) {#3};%
    \node[anchor=north, font=\ttfamily\large] at (71,-0.5) {#4};%
  \end{tikzpicture}%
}

% --- Honig-EAN: 4 270004 371635 ----------------------------------
% Pattern für Erstziffer 4: LGLLGG
% Linke Ziffern: 2(L), 7(G), 0(L), 0(L), 0(G), 4(G)
% Rechte Ziffern: 3(R), 7(R), 1(R), 6(R), 3(R), 5(R)
\newcommand{\barcodehonig}{%
  \eanbarcodebits
    {%
      1,0,1,%
      0,0,1,0,0,1,1, 0,0,1,0,0,0,1, 0,0,0,1,1,0,1,%
      0,0,0,1,1,0,1, 0,1,0,0,1,1,1, 0,0,1,1,1,0,1,%
      0,1,0,1,0,%
      1,0,0,0,0,1,0, 1,0,0,0,1,0,0, 1,1,0,0,1,1,0,%
      1,0,1,0,0,0,0, 1,0,0,0,0,1,0, 1,0,0,1,1,1,0,%
      1,0,1%
    }%
    {4}%
    {270004}%
    {371635}%
}
